\section{Final model-selection rationale and test-set evaluation}
\label{sec:final}

\subsection{Q5d: Multi-criteria selection}
Criteria (direction): performance = repeated-CV mean Macro-F1 ($\uparrow$); stability = repeated-CV SD
($\downarrow$); robustness = absolute validation Macro-F1 at the highest Gaussian-noise severity
($\sigma = 1.0$) ($\uparrow$); interpretability and implementation simplicity = ordinal 1--5 scores set in
\texttt{src/config.py} ($\uparrow$); training cost = fit time on the training partition ($\downarrow$);
inference cost = prediction time per 1000 rows ($\downarrow$). The two cost criteria are $\log_{10}$-transformed
before normalisation. Each criterion is then min-max normalised over the candidates so that 1 = best. The
normalised areas under both degradation curves are reported alongside (Table~\ref{tab:q5_robustness_criteria})
but are not used in the weighted score. Sensitivity analysis: fixed alternative schemes, leave-one-criterion-out,
and 5{,}000 Dirichlet weight vectors drawn (i) uniformly over the simplex ($\alpha = 1$) and (ii) centred on
the base weights ($\alpha = 20\,w$).
\inputtable{q5_robustness_criteria}
\inputtable{q5_cost}
\inputtable{q5_mcda_weights}
\inputtable{q5_mcda}
\inputtable{q5_sensitivity_schemes}
\inputtable{q5_dirichlet_wins}
\fig{q5_sensitivity_heatmap}{Model ranks under alternative weighting schemes.}

Our base MCDA weights reflect realistic retail-banking deployment priorities: predictive performance receives the largest share (0.40) to drive campaign ROI, followed by corruption robustness (0.20) to handle noisy CRM records, interpretability (0.15) to satisfy regulatory auditing and explain targeting decisions to campaign managers, CV stability (0.10) for operational reliability, and training cost, inference cost, and implementation simplicity (0.05 each). Naive Bayes receives a 4/5 interpretability score (vs.\ 2/5 for RF) because every prediction decomposes directly into additive log-odds contributions per feature, whereas RF requires post-hoc approximations (e.g., SHAP) across hundreds of deep trees. Under the base scheme, Naive Bayes ranks first (0.767), ahead of Random Forest (0.726) and Decision Tree (0.652). Sensitivity analysis shows that NB wins under base, equal, robustness-heavy, and cost-heavy weights as well as six of seven leave-one-criterion-out variants; however, RF takes first place under performance-heavy (0.60), performance-only, and \emph{base without interpretability} schemes (while DT wins under interpretability-heavy). In 5,000 centered Dirichlet weight perturbations ($\alpha = 20w$), NB wins 63.3\% of draws versus 32.8\% for RF (64.8\% vs.\ 3.7\% under uniform Dirichlet), indicating moderate robustness with clear sensitivity to how heavily interpretability is valued over raw Macro-F1. Methodologically, min--max normalization is vulnerable to outlier compression---for instance, DT's extreme CV SD (0.023) compresses all other models near 1.0, and stretching tiny runtime differences (e.g., 3.6\,ms vs.\ 9.4\,ms per 1,000 rows) onto $[0,1]$ can overstate negligible cost gaps despite $\log_{10}$ scaling---and because robustness ($\sigma=1.0$ Macro-F1) correlates with clean performance, predictive accuracy is partially double-counted.

\subsection{Q5e: Frozen configuration and single test evaluation}
Frozen configurations and the selection rule (\texttt{results/models/selection\_rule.json}) were written by
\texttt{src/q5\_robustness\_selection.py} before \texttt{src/q5\_final\_test.py} was executed once. Candidates were
refit on training+validation and evaluated once on the test partition; the run is logged in
\texttt{results/logs/final\_test\_runs.log}.
\inputtable{q5_test_results}
\inputtable{q5_test_per_class}
\fig[0.4]{q5_test_confusion_selected}{Test confusion matrix of the selected model.}
\inputtable{q5_test_results_regression}

After freezing the selection rule and refitting all models on the combined training and validation splits (85\% of data), the locked test set ($n = 6{,}782$) was evaluated once. The selected Naive Bayes model achieves a test Macro-F1 of 0.667 (95\% CI $[0.651, 0.684]$), accuracy of 0.857, ROC-AUC of 0.763, and PR-AUC of 0.391 (TN $= 5{,}468$, FP $= 521$, FN $= 448$, TP $= 345$). Its progression across CV (0.665), validation (0.662), and test (0.667)---mirrored by RF (0.697 $\to$ 0.706 $\to$ 0.704) and GB (0.664 $\to$ 0.674 $\to$ 0.666)---demonstrates near-zero generalization gap and validates the leakage-free protocol. Consistent with validation, Random Forest leads in raw test Macro-F1 at 0.704 ($[0.687, 0.720]$), with non-overlapping confidence intervals relative to NB; nevertheless, switching to RF after observing test outcomes would constitute test-set leakage and violate pre-registered MCDA selection, since NB was deliberately chosen for its multi-criteria balance rather than peak Macro-F1 alone. Decision Tree scores notably higher on validation (0.696) and test (0.699) than its CV mean (0.664), reflecting both its high fold-to-fold variance (SD $= 0.023$) and the benefit of refitting on 15\% additional data. Similarly, in the California Housing test evaluation, refitting on train+val helps lift test $R^2$ slightly above validation across all models (e.g., RF $0.807 \to 0.826$, DT $0.700 \to 0.742$, OLS $0.587 \to 0.604$), though all earlier dataset and protocol limitations remain applicable.
